\sectionguard
\section*{Appendix: Scope and Qualifications}

\textbf{As of 25 July 2026.} Everything below states the boundaries of this
guide as at that date. Retrieval mechanics, checksums, query detail, discarded
leads, and the full reception and proposal audits are in
\texttt{propers/verified.md} and \texttt{research/scope.md} beside this
document's source; they are not repeated here.

\editionnote

\textbf{Edition and formulary identity.} The work is the temporal formulary
printed \latin{DOMINICA NONA post Pentecosten}, \latin{II classis}, on pp.\
388--389 of the 1962 \latin{editio typica} of the \latin{Missale Romanum}
(Vatican City: Typis Polyglottis Vaticanis), with marginal numbers 1522--1531.
Green follows the general rubrics for Sundays after Pentecost and is not
printed on those pages. The formulary appoints the \latin{Credo} and the
Preface of the Most Holy Trinity, and appoints no Tract, Sequence, second
oration, blessing or ritual text.

\textbf{Text verification.} All ten elements, their headings, references,
wording, accents, punctuation, rank, rubrics and both formulary boundaries
were visually collated on 25 July 2026 against the CMAA facsimile of the
Vatican typical edition at 300 and 600 dpi, and independently against page
images of the 1962 Benziger \latin{editio iuxta typicam}. The two editions
differ at the points tabulated in the commentary; those are edition
differences, not doubtful readings, and nothing was blended or substituted.
The appointed text was separately collated against the Clementine Vulgate, and
the ten places where it diverges are tabulated in the commentary as verified
textual observations. Not printing the Latin in full changed what reaches the
page and changed nothing about that collation. The Internet Archive OCR was
used only to locate the formulary and controls no published wording.

\textbf{What this guide is not.} It is a hand-missal companion, not a
liturgical book, a critical edition, a homily, an assembly sheet for one civil
date, or a substitute for the sources cited. It does not treat occurrence,
commemorations, particular calendars, or celebration-mode branches for any
year; those belong to the edition-specific assembly workflow and require their
own evidence.

\textbf{Translation and received text.} The project has composed, translated,
paraphrased, modernised and adapted nothing. English for the scriptural
propers --- Introit, Gradual, Alleluia, Epistle, Gospel, Offertory, Communion
--- is the Douay--Rheims as revised by Bishop Challoner (1749--1752), the
English of the Clementine Vulgate the missal prints, quoted from the Project
Gutenberg e-text of that revision. English for the three orations is the
anonymous rendering in \work{The Roman Missal translated into the English
language for the use of the laity}, first revised edition (Philadelphia:
Eugene Cummiskey, 1861), taken from the three lines that book supplies for
this formulary and reproduced with its own spelling and abbreviated
conclusions. Both are in the public domain in the United States; neither is an
official liturgical translation, and neither was composed for the 1962 books.

\textbf{Four places where the English does not answer the Latin, and is not
made to.} The Introit's closing vocative \latin{protéctor meus, Dómine} has no
counterpart in Psalm 53 and none in the Douay; the Introit's psalm verse sings
\latin{líbera me} where the Clementine, and therefore the Douay, reads
\latin{judica me}, ``judge me''; the Communion's \latin{dicit Dóminus} is not
part of John 6:57; and the Offertory is a centonisation of four separated
verse-parts, so its four Douay verses are printed whole with the fragments
marked rather than stitched into a seamless paragraph. In each case the gap is
stated where it falls. No rendering of the project's own has been supplied
anywhere.

\textbf{Consulted for reception, not for English.} Dom Laurence Shepherd's
translation of Guéranger, \work{The Liturgical Year} \dots\ \work{Time after
Pentecost}, vol.~II, 2nd ed.\ (Stanbrook Abbey: Burns \& Oates, 1909),
supplied the English of the appointed texts in an earlier state of this guide
and no longer does. It is retained only as a witness to the devotional
reception of this formulary, including the supersessionist reading discussed
and rejected in the detailed commentary; citing it there is not endorsement of
its exposition.

\textbf{Search bounds.} Reception was searched in Latin, Greek in
translation, and English, across the major reasonably accessible patristic,
medieval and later corpora relevant to each appointed passage: the
\latin{Enarrationes}, Tractates, sermons, letters and principal treatises of
Augustine; Ambrose on Luke; Origen's Lucan homilies in Jerome's Latin;
Gregory's Gospel homilies; Bede on Luke; Chrysostom on First Corinthians and
John; Cyril of Alexandria on Luke and John; Ambrosiaster; Theodoret on the
Psalms; Aquinas's psalm, Johannine and Pauline commentaries, the
\work{Catena aurea} and the complete \work{Summa theologiae}; and Bellarmine
on the Psalms. Consequential negative results --- notably that Augustine has
no exegesis of Luke 19:41--44, that the \work{Summa} never cites this Gospel,
that Aquinas wrote no commentary on Psalm 58, and that the exposition of
1~Corinthians 10 in the Thomistic corpus is Peter of Tarantaise's --- are
reported in the commentary rather than concealed. No exhaustive search of
Syriac originals, chant manuscripts, untranslated homiliaries, subscription
databases or current specialist monographs is claimed, and none of these
negatives is offered as a claim about literature outside the corpora named.

\textbf{Bounds carried by particular claims.} Theodoret was consulted in a
modern English translation that is in copyright, so he is summarised and never
quoted, and the PG column numbers given for him are those printed inline in
that translation rather than read from a Migne scan. Bede's PL 92 columns come
from a transcription's embedded markers, not from a facsimile. Column numbers
for Chrysostom, Cyril and Augustine's Tractates were not verified beyond the
volume level and are therefore not printed. Johner's 1940 chant commentary
shows no copyright renewal in the standard record, which points to but does
not establish public-domain status; it is paraphrased and quoted only in short
compass. Origen's homily has no public-domain English, so only its Latin is
quoted. Augustine's second sermon on Psalm 18 is not in the public-domain
English collection and was read in Latin.

\textbf{Psalm numbering.} Three series are in use and two of them are
routinely both called ``modern,'' so the dossier sheet prints the two that
differ and this appendix supplies the third. \textit{Vulg.} is the
Vulgate/Septuagint number, which the missal cites and which the Douay--Rheims
carries throughout. \textit{Eng.} is the convention of the King James and
Revised Standard versions, which leave the psalm inscription unnumbered and so
run one or two verses behind the Vulgate; the dossiers print it beside the
Vulgate number. \textit{Heb.} is the Masoretic number of the Nova Vulgata and
the New American Bible, which renumbers the psalm but keeps the Vulgate's
verse division: Ps.\ 8:2 Vulg.\ is Heb.\ 8:2; Ps.\ 18:9--12 is Heb.\
19:9--12; Ps.\ 53:3, 6--7 is Heb.\ 54:3, 6--7; Ps.\ 58:2 is Heb.\ 59:2. The
conversions were taken from the concordance registered with the Douay--Rheims
edition used here, not computed by this guide; all four psalms of this
formulary have a uniform offset, so it holds from the head of each psalm.

\textbf{Historical bounds.} Dates and places in the page-2 dossiers
distinguish traditional attribution from current historical judgment and are
orienting rather than probative. This guide does not settle the year of the
Passion, the date or place of Luke's or John's composition, or the historical
authorship of any psalm.

\textbf{Rights.} No public-domain conclusion is claimed for the 1962 typical
edition. This repository's rights record for the facsimile finds that edition's
own status unresolved, and more likely than not restored in the United States
under 17 U.S.C.\ \S104A to the end of 2057 (that record is
\texttt{missale-romanum-1962-facsimile-rights-v1.toml}, under
\texttt{src/sources/inventories/}).
What permits the Latin quoted here is 17 U.S.C.\ \S103(b): the copyright in a
revised edition extends only to what that edition itself contributed, and ``does
not affect or enlarge the scope, duration, ownership, or subsistence of'' any
copyright in the preexisting material. This formulary is preexisting material.
All ten appointed elements of it --- Introit, Collect, Epistle, Gradual,
Alleluia, Gospel, Offertory, Secret, Communion and Postcommunion, in the same
order and at the same scriptural references --- are printed in the Pustet
printing of Ratisbon, 1862, which this repository tracks as public-domain text,
and six of the ten in the tracked Venice Missal of 1570 besides, so the 1962
edition's copyright does not reach them; the element-by-element check and its
bounds are recorded in \texttt{propers/verified.md}.
That is this project's own reading of the statute and of the evidence,
not a clearance anyone granted, and it is not legal advice. No page image of
either digitisation is reproduced here, and none of the 1962 edition's own new
matter --- the 1955 Holy Week, the 1960 Code of Rubrics, the sanctoral added
after 1920, its front matter --- is published in this guide; the edition's rank,
pagination and marginal numbers are cited as facts about the book. The project
claims no rights in the received liturgical text. Both digitisations used to
verify it are third-party artifacts whose own
rights status the repository's source records list as unresolved; neither is
redistributed here. Quotations from sources still in copyright are avoided or
kept to short compass with attribution, as noted at each point of use. The
common reuse notice on the final page states the outbound boundary and does
not relicense any of this material. Nothing here is permission to extract and
redistribute the incorporated liturgical text independently.

\textbf{Review state.} This document has had internal production review only:
source collation, build, and page-by-page visual inspection by the same agent
that produced it. It has received no independent liturgical, theological,
specialist, rights or ecclesiastical review, and carries no imprimatur, nihil
obstat or ecclesiastical approval. The exact published PDF snapshots have
distribution authorization under the release record; that authorization does
not relicense the incorporated liturgical text for independent extraction and
does not constitute any of the outstanding reviews.
